% This latex document tries to copy the microsoft word template of XXXIV CILAMCE

\documentclass[a4paper,12pt,fleqn]{article}

\usepackage{graphicx}
\usepackage{amsmath,amsfonts, amsthm}
\usepackage{enumerate}
\usepackage[shortlabels]{enumitem} 
\usepackage{tabularx}
\usepackage{comment}
\usepackage{url}
\usepackage{color}
\usepackage{cancel}
\usepackage{subfigure}


% PACKAGES USED - packages that need to be previously installed on your computer
\usepackage[lmargin=2.5cm, rmargin=2.5cm, tmargin=2.5cm, bmargin=3.5cm ]{geometry}
\usepackage{graphicx}
%\usepackage{times}
\usepackage{indentfirst}
\usepackage{fancyhdr}
%\usepackage{titlesec}
%\usepackage[english]{babel}      % nomes em português
%\selectlanguage{portuges}
\usepackage[english,brazil]{babel} %Opções de escrita em inglês, português para Figura, Tabela.
\usepackage{parskip}
\usepackage{setspace}

%%% Overleaf:
\usepackage[utf8]{inputenc}
%%% Texmaker:
% \usepackage[latin1]{inputenc}
\usepackage[brazil]{babel}


\usepackage{amssymb} % Usado para desigualdades nas equações
\usepackage{dsfont} %Usado para conjuntos N, Z, Q, R, C


\usepackage{cite}


\usepackage[pdftex]{hyperref}    % "links" dentro do pdf
\usepackage[T1]{fontenc}          % permitir \hyphenation

\usepackage{pifont}				% caracteres decorativos
%\usepackage[utf8]{inputenc} % Usado para acentuações



% CONFIGURATION
\renewcommand*\arraystretch{1.5} % height of the table's row
\usepackage[font=bf, size=footnotesize]{caption} % bold face font for labels of figures and tables
\renewcommand*\thesection{\arabic{section}} % for use book class instead article class
%\hyphenpenalty=10000 % uncomment for avoir hypenization
%\titleformat*{\section}{\large\bfseries}
%\titleformat*{\subsection}{\large\bfseries}
%\titlespacing\section{0pt}{12pt plus 12pt minus 6pt}{6pt plus 6pt minus 6pt}
%\titlespacing\subsection{0pt}{6pt plus 6pt minus 6pt}{6pt plus 6pt minus 6pt}
%\titlespacing\subsubsection{0pt}{6pt plus 6pt minus 6pt}{6pt plus 6pt minus 6pt}
\setlength{\parskip}{6pt}
\setlength{\parindent}{0.75cm}
\let\ds=\displaystyle
% PDF zoom {the last parameter is the zoom: 1.00 (100%), 0.5 (50%), ...}
\usepackage{hyperref}
\hypersetup{pdfstartview={XYZ null null 1.00}}

\usepackage{lipsum}

\usepackage{bigints}


%%% Para usar sen em portugues
\DeclareMathOperator{\sen}{sen}
\DeclareMathOperator{\arctg}{arctg}

\usepackage[table]{xcolor}


%%% Para graficos:
\usepackage{tikz}
\usepackage{pgfplots}


\begin{document}

{\color{blue}

\begin{center}
\large
Cálculo Numérico - IME/UERJ 

Gabarito - Lista de Exercícios 2 - Raízes de funções
\end{center}

}

\vspace{0.8cm}

  \begin{enumerate}
    
     
    \item 
    
    {\color{blue}
    
    \textbf{Resposta:} 
    
%    Vamos calcular os valores de $p(x)$ em pontos suficientemente próximos das raízes. Assim, temos a tabela:
%
%		\begin{center}
%  				\begin{tabular}{ | l | l | l | l | l | l | l | l | l | l | l |}
%   				 \hline
%   					 $x$ & 0 & 0,5 & 1 & 1,5 & \cellcolor{yellow}2 & \cellcolor{yellow} 2,5 & \cellcolor{yellow} 3 & 3,5 & 4 & 4,5 \\ \hline
%   					 $p(x)$ & 400 & 85,75 & 0 & -7,8125 & \cellcolor{yellow} - 2 & \cellcolor{yellow} 0 & \cellcolor{yellow} - 0,5 & - 0,3125 & 0 & 1,75 \\ \hline
%				\end{tabular}
%		\end{center}
%		
%		Analisando o sinal de $p(x)$ para os valores da tabela, concluímos que nas proximidades da raiz $2,5$, no intervalo $(2,3)$, a função $p(x)$ não mudou de sinal. Portanto, o método da Bisseção não funciona para a raiz $2,5$.

2,5
    
    } 
    
    \vspace{0.3cm}
    
    \item 
    
    {\color{blue}
    
    $[2,3]$. Número de iterações: $k_{\text{min}} = 10$.
    
    }   
    
    \vspace{0.6cm}    
      

\item 

{\color{blue}


\begin{enumerate}[(a)]

\item II. $\vert g_b'(1,5) \vert < 1$ (Verifique!)
\item 11

\end{enumerate}

}

\vspace{0.3cm}  

    \item 

    {\color{blue}
    
   Como $\vert \varphi_1'(2) \vert < \vert \varphi_2'(2) \vert < 1$, logo $\varphi_1(x)$ gera sequências mais rapidamente convergentes para a raiz.
    
    }   
    
\vspace{0.3cm}   


\item

\begin{enumerate}[(a)]
\item

      {\color{blue} 
      
         \textbf{Resposta:}
         
%         $f(x) = e^{-2x} + x^2 - 4 = 0 \Rightarrow e^{-2x} = 4 - x^2$.
         
%         Portanto, graficamente, as raízes são os pontos de interseção das funções $e^{-2x}$ e $9 - x^2$.
     
%\begin{tikzpicture}
%\begin{axis}[
%    axis lines = center,
%    xtick distance = 1,
%%    ytick distance = 0.25,
%    ymin=0, ymax=8,  
%    xmin=-3, xmax=3,    
%    xlabel = \(x\),
%    ylabel = {\( y \)},
%]
%%Below the red parabola is defined
%\addplot [
%    domain=-3:3, 
%    samples=100, 
%    color=red,
%]
%{exp(-2*x)};
%\addlegendentry{\( e^{-2x} \)}
%%Here the blue parabola is defined
%\addplot [
%    domain=-4:4, 
%    samples=100, 
%    color=black,
%    ]
%    {4 - x^2};
%\addlegendentry{\( 4 - x^2 \)}
%
%\addplot[mark=*] coordinates { (-0.639263, 3.59134) };
%\addplot[mark=*] coordinates { (1.995375, 0.01848) };
%
%\end{axis}
%\end{tikzpicture}

Usando gráficos ou o Teorema do Valor Intermediário, temos duas raízes: $r_1 \in (-\ln(4)/2,0) \approx (-0,69;0)$ e $r_2 \in (1,2)$.

\textbf{Obs.:} O denominador de $\varphi_1'(x)$ é $\sqrt{4-e^{-2x}}$. 

Logo, $4-e^{-2x} > 0 \Rightarrow x > -\ln(4)/2 \approx - 0,69$.

Portanto, para $r_1$, uma boa aproximação inicial é $x_0 = -0,5$, enquanto para $r_2$, uma boa aproximação inicial é $x_0 = 1,9$.     
     
     
       }
        
        
        \vspace{0.3cm}  
        
        \item

         {\color{blue} 
      
         \textbf{Resposta:}
         
%         Pelo teorema do método do ponto fixo, há uma sequência convergente para uma raiz quando $\varphi'(x) \leq M < 1$ para todo $x \in I$, onde $I$ é um intervalo centrado na raiz e $x_0 \in I$.
%         
%         $\vert \varphi_1'(x) \vert = \Big \vert \dfrac{e^{-2x}}{\sqrt{4 - e^{-2x} }} \Big \vert = \dfrac{e^{-2x}}{\sqrt{4 - e^{-2x} }} $
%         
%         Então, para as aproximações iniciais das raízes do item (a):
         
         $\vert \varphi_1'(r_1 \approx -0,5) \vert \approx {\color{red} 2,4010 > 1 \Rightarrow \text{Diverge!}} $
         
         $\vert \varphi_1'(r_2 \approx 1,9) \vert \approx 0,0112 < 1 \Rightarrow \varphi_1(x) \text{ converge para } r_2 \in (1,2)$.
         
         \vspace{0.3cm}
         
%        $\vert \varphi_2'(x) \vert = \Big \vert  \dfrac{x}{4-x^2} \Big \vert$
%         
%         Então, para as aproximações iniciais das raízes do item (a):
         
         $\vert \varphi_2'(r_1 \approx -0,5) \vert \approx 0,1333 < 1 \Rightarrow \varphi_2(x) \text{ converge para } r_1 \in (-1,0)$.
         
         $\vert \varphi_2'(r_2 \approx 1,9) \vert \approx {\color{red} 4,8718 > 1 \Rightarrow \text{Diverge!} } $
         
         \item
         
         Usando $x_0 = -0,5$ para a raiz negativa e $x_0 = 1,9$ para raiz positiva, os resultados para o método de Newton-Raphson com tolerância $\epsilon \leq 10^{-4}$ são:
         
         Raiz negativa: $r_1 \approx - 0,6393$;
         Raiz positiva: $r_2 \approx 1,9954$.
       
         }                
        
\end{enumerate}

\vspace{0.3cm}
        
    
%\item 
%
%{\color{blue} 
%
%A maior raiz é $r_2 \approx 1,9954$.
%     
%}   

%%\item 
%%
%%{\color{blue} 
%%
%%\begin{enumerate}
%%
%%\item 
%%
%%Usando gráficos ou Teorema do Valor Intermediário, escolha uma boa aproximação inicial no intervalo $(1,2)$. Um exemplo é $x_0 = 1,5$.
%%
%%\item
%%
%%$\vert \varphi'_1(1,5) \vert \approx 1,5493 > 1$ (não converge), $\vert \varphi'_2(1,5) \vert \approx 0,2171 < 1$ (converge). 
%%
%%\item
%%
%%$r \approx 1,5569$
%%
%%\end{enumerate}
%%
%%}

%\vspace{0.3cm}
%
%
%    \item 
%      \label{exer12}
%      
%    {\color{blue}
%
%    
%      \begin{enumerate}[(a)]
%        \item 
%        
%        $\varphi_1(x) = e^{-2x}$. Como $\vert \varphi_1'(x) \vert = 2 \cdot e^{-2x} < 1 \Rightarrow x > 0.34657$, um bom intervalo é $(0.34657,0.5)$, pois $f_1(0.34657) \cdot f_1(0.5) < 0$.
%               
%        \item 
%        
%%        Vejamos as estimativas para as raízes graficamente:
%%        
%%        $f_2(x) = 0 \Rightarrow \ln(x) = x-2$. 
%%        
%%\begin{tikzpicture}
%%        
%%        
%%\begin{axis}[
%%    axis lines = center,
%%    xtick distance = 0.5,
%%%    ytick distance = 0.25,
%%    xlabel = \(x\),
%%    ylabel = {\(  y \)},
%%    legend style={at={(0.2,1.0)},anchor=north west}
%%]
%%%Below the red parabola is defined
%%\addplot [
%%    domain=0:4, 
%%    samples=100, 
%%    color=red,
%%]
%%{ln(x)};
%%\addlegendentry{\( \ln(x) \)}
%%%Here the blue parabola is defined
%%\addplot [
%%    domain=0:4, 
%%    samples=100, 
%%    color=black,
%%    ]
%%    { x-2 };
%%\addlegendentry{\( x-2 \)}
%%
%%\addplot[mark=*] coordinates { (0.1586, -1.8414) };
%%\addplot[mark=*] coordinates { (3.1462, 1.1462) };
%%
%%\end{axis}
%%\end{tikzpicture}
%
%As raízes são: $r_1 \in (0,0.5)$, $r_2 \in (3,3.5)$.
%
%Então, uma boa estimativa inicial de raiz para $r_1$ é $x_0 = 0.2$ e uma boa estimativa inicial para $r_2$ é $x_0 = 3.2$.
%
%%Fazendo $f_2(x) = 0$ e isolando a variável $x$, obtemos:
%%
%%$\ln(x) = x - 2 \Rightarrow x = e^{x-2}$.
%
%%Logo,
%Uma função de iteração é $\varphi_2(x) = e^{x-2}$.
%
%%Pelo Teorema do Método do Ponto Fixo, quando $\vert \varphi'_2(x) \vert < 1$ para todo $x \in I$, sendo $I$ centrado na raiz, existe uma sequência $\{x_k\}$ que converge para a raiz usando a iteração $x_{k+1} = \varphi_2(x_k)$.
%%
%%Então, vamos analisar o que acontece com as estimativas iniciais:
%%
%%Sabemos que $\vert \varphi'_2(x) \vert = e^{x-2}$.
%%
%%Assim, 
%
%Para $r_1$: $\vert \varphi'_2(0.2) \vert = e^{0.2-2} = e^{-1.8} \approx 0.1653 < 1 \Rightarrow \varphi_2(x)$ converge para $r_1$.
%
%Para $r_2$: $\vert \varphi'_2(3.2) \vert = e^{3.2-2} = e^{1.2} \approx {\color{red} 3.3201 > 1 \Rightarrow \varphi_2(x) }$ {\color{red} não converge para $r_2$}.
%
%%Então, neste item, para que tenhamos uma função de iteração que gere uma sequência convergente para a raiz, um intervalo $(a,b)$ para a raiz é $(0,0.5)$ e uma função de iteração associada é $\varphi_2(x) = e^{x-2}$.
%
%        
%        \item $\varphi_3(x) = e^{x/6}$; $r \in (1,2)$.
%        \item $\varphi_{41}(x) = \sqrt{\sen(x)}$; $r_1 \in (0.8, 0.9)$. A  raiz $r_2 \in (-0.1,0.1)$ deve ser encontrada em $\varphi_{42}(x) =  - \sqrt{\sen(x)}$.
%        
%        \item 
%        
%%       Vejamos as estimativas para as raízes graficamente:
%%       
%%       $f_5(x) = 0 \Rightarrow \dfrac{x}{4} = \cos(x)$
%%        
%%        \begin{tikzpicture}
%%\begin{axis}[
%%    axis lines = center,
%%    xtick distance = 1,
%%%    ytick distance = 0.25,
%%    ymin=-2, ymax=3,  
%%    xmin=-5, xmax=2, 
%%    xlabel = \(x\),
%%    ylabel = {\(  y \)},
%%    legend style={at={(0.2,1.0)},anchor=north west}
%%]
%%
%%%\addplot [
%%%    domain=0:4, 
%%%    samples=100, 
%%%    color=red,
%%%]
%%%{rad(acos(x/4))};
%%%\addlegendentry{\( arc \cos(x/4) \)}
%%
%%\addplot [
%%    domain=-5:2, 
%%    samples=100, 
%%    color=red,
%%    ]
%%    { cos(deg(x)) };
%%\addlegendentry{\( \cos(x) \)}
%%
%%\addplot [
%%    domain=-5:2, 
%%    samples=100, 
%%    color=black,
%%]
%%{ x/4 };
%%\addlegendentry{\( x/4 \)}
%%
%%\addplot[mark=*] coordinates { (1.2524, 0.3130) };
%%\addplot[mark=*] coordinates { (-2.133, -0.5331) };
%%\addplot[mark=*] coordinates { (-3.595, -0.8990) };
%%
%%
%%
%%\end{axis}
%%\end{tikzpicture}
%
%%Vemos que 
%As raízes são: $r_1 \in (-4,-3)$, $r_2 \in (-3,-2)$ e $r_3 \in (1,\pi/2)$.
%
%Então, uma boa estimativa inicial de raiz para $r_1$ é $x_0 = -3.8$, uma boa estimativa inicial para $r_2$ é $x_0 = -2.1$ e uma boa estimativa inicial para $r_3$ é $x_0 = 1.2$.
%
%%Fazendo $f_5(x) = 0$ e isolando a variável $x$, obtemos:
%%
%%$\dfrac{x}{4} = \cos(x) \Rightarrow x = \arccos(x/4)$.
%
%%Logo, 
%Uma função de iteração que podemos usar é $\varphi_5(x) = \arccos(x/4)$.
%
%Então, vamos analisar o que acontece com as estimativas iniciais:
%
%%Sabemos que $\vert \varphi'_5(x) \Big\vert = \vert - \dfrac{1}{4 \sqrt{1 - (x/4)^2}} \Big\vert = \dfrac{1}{4 \sqrt{1 - (x/4)^2}} = \dfrac{1}{\sqrt{16 - x^2}}$.
%%
%%Para $r_1$: $\vert \varphi'_5(-3.8) \vert = \dfrac{1}{\sqrt{16 - (-3.8)^2}} \approx 0.8006 < 1 \Rightarrow \varphi_5(x)$ converge para $r_1$.
%%
%%Para $r_2$: $\vert \varphi'_5(-2.1) \vert = \dfrac{1}{\sqrt{16 - (-2.1)^2}} \approx 0.2937 < 1 \Rightarrow \varphi_5(x)$ converge para $r_2$.
%%
%%Para $r_3$: $\vert \varphi'_5(1.2) \vert = \dfrac{1}{\sqrt{16 - (1.2)^2}} \approx 0.2621 < 1 \Rightarrow \varphi_5(x)$ converge para $r_3$.
%
%%Sabemos que $\vert \varphi'_5(x) \Big\vert = \vert - \dfrac{1}{4 \sqrt{1 - (x/4)^2}} \Big\vert = \dfrac{1}{4 \sqrt{1 - (x/4)^2}} = \dfrac{1}{\sqrt{16 - x^2}}$.
%
%Para $r_1$: $\vert \varphi'_5(-3.8) \vert \approx 0.8006 < 1 \Rightarrow \varphi_5(x)$ converge para $r_1$.
%
%Para $r_2$: $\vert \varphi'_5(-2.1) \vert \approx 0.2937 < 1 \Rightarrow \varphi_5(x)$ converge para $r_2$.
%
%Para $r_3$: $\vert \varphi'_5(1.2) \vert \approx 0.2621 < 1 \Rightarrow \varphi_5(x)$ converge para $r_3$.
%
%%
%%Então, neste item, para que tenhamos uma função de iteração que gere uma sequência convergente para a raiz, temos três intervalos $(-4,-3)$, $(-3,-2)$ e $(1,\pi/2)$, respectivamente, para as raízes $r_1$, $r_2$, $r_3$ e uma função de iteração associada $\varphi_5(x) = \arccos(x/4)$.
        
        
%      \end{enumerate} 

%    }   
%    
%    \vspace{0.6cm}       
    



   \item {\color{blue} Resolver. }

\end{enumerate}


\end{document}
